\documentclass[a4paper]{article}
\addtolength{\hoffset}{-2.25cm}
\addtolength{\textwidth}{4.5cm}
\addtolength{\voffset}{-3.25cm}
\addtolength{\textheight}{5cm}
\setlength{\parskip}{0pt}
\setlength{\parindent}{0in}

\usepackage[square,sort,comma,numbers]{natbib}
\usepackage{blindtext} % Package to generate dummy text
\usepackage{charter} % Use the Charter font
\usepackage[utf8]{inputenc} % Use UTF-8 encoding
\usepackage{microtype} % Slightly tweak font spacing for aesthetics
\usepackage{amsthm, amsmath, amssymb} % Mathematical typesetting
\usepackage{float} % Improved interface for floating objects
\usepackage{hyperref} % For hyperlinks in the PDF
\usepackage{graphicx, multicol} % Enhanced support for graphics
\usepackage{xcolor} % Driver-independent color extensions
\usepackage{pseudocode} % Environment for specifying algorithms in a natural way
\usepackage[mmddyy]{datetime} % Uses YEAR-MONTH-DAY format for dates

\usepackage{fancyhdr} % Headers and footers
\pagestyle{fancy} % All pages have headers and footers
\fancyhead{}\renewcommand{\headrulewidth}{0pt} % Blank out the default header
\fancyfoot[L]{} % Custom footer text
\fancyfoot[C]{} % Custom footer text
\fancyfoot[R]{\thepage} % Custom footer text
\newcommand{\note}[1]{\marginpar{\scriptsize \textcolor{red}{#1}}} % Enables comments in red on margin

\DeclareMathOperator*{\argmin}{arg\,min}

%----------------------------------------------------------------------------------------


%-------------------------------
%	TITLE VARIABLES (identify your work!)
%-------------------------------

\newcommand{\yourname}{Nathan Ugwonali} % replace YOURNAME with your name
\newcommand{\youremail}{nathan.ugwonali@ens-paris-saclay.fr} % replace YOUREMAIL with your email
\newcommand{\assignmentnumber}{5} % replace X with the lab session number

\begin{document}

%-------------------------------
%	TITLE SECTION (do not modify unless you really need to)
%-------------------------------
\fancyhead[C]{}
\hrule \medskip
\begin{minipage}{0.295\textwidth} 
\raggedright
\footnotesize
\yourname \hfill\\
\youremail
\end{minipage}
\begin{minipage}{0.4\textwidth} 
\centering 
\large 
Lab session \# \assignmentnumber\\ 
\normalsize 
ALTEGRAD 2023\\ 
\end{minipage}
\begin{minipage}{0.295\textwidth} 
\raggedleft
\today\hfill\\
\end{minipage}
\medskip\hrule 
\bigskip



%-------------------------------
%	ASSIGNMENT CONTENT (add your responses)
%-------------------------------

\section{Question 1}

$z_1^{(2)} = z_4^{(2)}$.

\paragraph{Derivation.} Let $W$, $a$ be the weights of the second layer and write the raw score and the update as
\begin{equation}
    e_{ij} = \text{LeakyReLU}\big(a^\top [W z_i^{(1)} \| W z_j^{(1)}]\big), \qquad
    \alpha_{ij} = \frac{\exp(e_{ij})}{\sum_{k \in \mathcal{N}(v_i)} \exp(e_{ik})}, \qquad z_i^{(2)} = \sum_{j \in \mathcal{N}(v_i)} \alpha_{ij} W z_j^{(1)} .
\end{equation}
For $v_1$, with $\mathcal{N}(v_1) = \{v_2, v_3\}$,
\begin{equation}
    \alpha_{12} = \frac{e^{e_{12}}}{e^{e_{12}} + e^{e_{13}}}, \qquad
    \alpha_{13} = \frac{e^{e_{13}}}{e^{e_{12}} + e^{e_{13}}}, \qquad
    z_1^{(2)} = \alpha_{12} W z_2^{(1)} + \alpha_{13} W z_3^{(1)} .
\end{equation}
For $v_4$, with $\mathcal{N}(v_4) = \{v_2, v_3, v_5, v_6\}$: each $e_{ij}$ depends only on $z_i^{(1)}$ and $z_j^{(1)}$, so $z_4^{(1)} = z_1^{(1)}$, $z_6^{(1)} = z_2^{(1)}$ and $z_5^{(1)} = z_3^{(1)}$ give
\begin{equation}
    e_{42} = e_{46} = e_{12}, \qquad e_{43} = e_{45} = e_{13} .
\end{equation}
The softmax denominator of $v_4$ is therefore $2\big(e^{e_{12}} + e^{e_{13}}\big)$, twice that of $v_1$, hence
\begin{equation}
    \alpha_{42} = \alpha_{46} = \frac{\alpha_{12}}{2}, \qquad \alpha_{43} = \alpha_{45} = \frac{\alpha_{13}}{2},
\end{equation}
and
\begin{equation}
    z_4^{(2)} = \frac{\alpha_{12}}{2}\big(W z_2^{(1)} + W z_6^{(1)}\big) + \frac{\alpha_{13}}{2}\big(W z_3^{(1)} + W z_5^{(1)}\big)
    = \alpha_{12} W z_2^{(1)} + \alpha_{13} W z_3^{(1)} = z_1^{(2)} .
\end{equation}
If the layer is followed by an activation $f$ (as in the GNN of Section 3.2), the equality still holds since $f$ is applied to the same vector.

\paragraph{Why.} The attention coefficients of each node are a softmax over its neighbours and sum to $1$, so the layer computes a \emph{weighted mean} of the neighbours' messages. The neighbourhood of $v_4$, seen as the multiset of representations $\{z_2, z_2, z_3, z_3\}$, is that of $v_1$, $\{z_2, z_3\}$, with every element duplicated: the proportions are identical. Duplicating neighbours doubles each exponentiated score but also the denominator, so the weights halve and the aggregate is unchanged. The layer thus cannot distinguish $v_1$ from $v_4$ although $v_4$ has twice as many neighbours: degree information is lost. A sum aggregator, as in GIN~\cite{xu2019powerful}, would give $2\,z_1^{(2)}$ for $v_4$ in this setting and separate the two nodes.


\section{Question 2}

No. Suppose all rows of $X$ are equal to some vector $x$. In the first layer, $W x_i = W x$ for every node, so all scores $e_{ij}$ are equal and the softmax gives uniform weights $\alpha_{ij} = 1/|\mathcal{N}(v_i)|$. Every node then receives
\begin{equation}
    z_i^{(1)} = f\Big(\sum_{j \in \mathcal{N}(v_i)} \frac{1}{|\mathcal{N}(v_i)|}\, W x\Big) = f(W x),
\end{equation}
so the representations are still identical after the first layer and, by the same argument, after the second one and the fully-connected layer. Because the attention weights form a convex combination, not even the degree of a node is captured. All nodes receive the same output distribution and hence the same predicted class: the model reduces to a constant classifier, and its accuracy is the proportion of the predicted class in the test set, regardless of training.

We verified this empirically by replacing the random features with a matrix of ones in \texttt{gnn\_karate.py}: over 3 runs, the model predicts class 1 for all 7 test nodes (4 of class 0, 3 of class 1), giving a test accuracy of $3/7 \approx 0.43$, against $1.00$ with random features. Random features act as (approximately) unique node identifiers that break this symmetry and let the model use the graph structure to separate the two communities.


\section{Question 3}

\begin{center}
\begin{tabular}{c|c|c|c}
     & sum & mean & max \\
    \hline
    $z_{G_1}$ & $[2.9,\ 2.3,\ 1.9]$ & $[0.97,\ 0.77,\ 0.63]$ & $[2.2,\ 1.8,\ 1.5]$ \\
    $z_{G_2}$ & $[3.4,\ 1.9,\ 4.3]$ & $[0.85,\ 0.48,\ 1.08]$ & $[2.2,\ 1.8,\ 1.5]$ \\
    $z_{G_3}$ & $[1.8,\ 1.2,\ 1.6]$ & $[0.90,\ 0.60,\ 0.80]$ & $[2.2,\ 1.8,\ 1.5]$ \\
\end{tabular}
\end{center}

The max readout maps the three graphs to the same vector, so it cannot distinguish them at all: in each column, the maximum ($2.2$, $1.8$, $1.5$) is reached by a node of every graph. Mean and sum both yield three distinct vectors here, but the sum distinguishes the graphs best: its representations are the most spread apart, and it additionally encodes the number of nodes, which the mean discards. In general, the mean cannot separate a multiset of node representations from the same multiset with every element repeated (cf.\ Question~1), and the max only retains the set of distinct values, whereas the sum is injective on a much larger class of multisets, which makes it the most expressive of the three readouts~\cite{xu2019powerful}.


\section{Question 4}

$z_{G_2} = 2\, z_{G_1}$.

$C_4$ and $C_8$ are both 2-regular and all node features equal $1$, so every row of $\tilde{A} = A + I$ contains exactly three ones. We show by induction that, within and across both graphs, all nodes have the same representation after every layer. If all rows of $Z$ are equal to a row vector $r$, then each row of $\tilde{A} Z$ is $3r$, so all rows of $\tilde{A} Z W$ equal $3rW$, and applying the same element-wise ReLU keeps them identical. Starting from $X = \mathbf{1}$, every row of $Z^{(2)}$ is therefore the same vector $h$ in both graphs (the biases of the linear layers do not change this argument). The sum readout then gives
\begin{equation}
    z_{G_1} = \sum_{i=1}^{4} h = 4h, \qquad z_{G_2} = \sum_{i=1}^{8} h = 8h = 2\, z_{G_1}.
\end{equation}
We verified this numerically with our implementation (randomly initialised model, $C_4$ and $C_8$ in the same batch): the input of the first fully-connected layer satisfies $z_{G_2} = 2 z_{G_1}$.

The message passing layers cannot tell the two cycles apart locally: like the 1-WL test, they see the same degree everywhere, so every node of $C_4$ looks exactly like every node of $C_8$. Only the sum readout separates the two graphs, and only through their number of nodes. With a mean or max readout, we would get $z_{G_1} = z_{G_2}$, and the model would assign the same prediction to both graphs.


%------------------------------------------------

\bibliographystyle{plain}
\bibliography{references} % citation records are in the references.bib document

\end{document}
